% =====================================================================
%  2bat-ud1-9.tex  —  SESSIÓ 9: El producte pas a pas (mecanització)
%  (sense preàmbul: s'inclou des de main.tex)
%  Criteri: C3 (producte de matrices). Mecanització: bateria graduada.
% =====================================================================
\horatitol{Sessió 9 — El producte pas a pas (mecanització)}%
{Guanyar destresa amb una bateria graduada de productes, de matrices petites a rectangulars.}

\subsection{Idea clau}

Ja saps què significa el producte de matrices i per a què serveix. Aquesta sessió
té un únic objectiu: que el \textbf{mecanisme} fila-per-columna et surti de manera
segura i ràpida, sense errors, encara que el problema tingui números més grans o
matrices rectangulars.

\begin{idea}
\textbf{Recordatori del mecanisme (i els dos paranys)}
\begin{itemize}[leftmargin=1.4em, itemsep=2pt]
  \item Cada casella del resultat surt de creuar una \textbf{fila} de la primera
        matriu amb una \textbf{columna} de la segona: es multiplica element a
        element i se \emph{sumen} els productes.
  \item La casella de la \textbf{fila $i$, columna $j$} del resultat surt de la
        fila $i$ de la primera i la columna $j$ de la segona.
  \item \textbf{Parany 1:} equivocar-se en la correspondència fila-columna (creuar
        una fila amb la fila enlloc de la columna).
  \item \textbf{Parany 2:} perdre's en la suma de productes (oblidar un terme o
        sumar malament). Escriu sempre tots els productes abans de sumar.
\end{itemize}
\end{idea}

\subsection{Exemples}

\begin{exemple}[un producte $2\times 2$ recordant el mecanisme]
\[
\begin{pmatrix} 2 & 3 \\ 1 & 4 \end{pmatrix}\cdot\begin{pmatrix} 5 & 0 \\ 2 & 6 \end{pmatrix}
=\begin{pmatrix} 2\per 5+3\per 2 & 2\per 0+3\per 6 \\ 1\per 5+4\per 2 & 1\per 0+4\per 6 \end{pmatrix}
=\begin{pmatrix} 16 & 18 \\ 13 & 24 \end{pmatrix}.
\]
\end{exemple}

\begin{exemple}[un producte rectangular]
Una matriu $2\times 3$ per una $3\times 2$ dona una matriu $2\times 2$:
\[
\begin{pmatrix} 1 & 2 & 0 \\ 3 & 1 & 2 \end{pmatrix}\cdot\begin{pmatrix} 2 & 1 \\ 0 & 4 \\ 3 & 1 \end{pmatrix}
=\begin{pmatrix} 1\per 2+2\per 0+0\per 3 & 1\per 1+2\per 4+0\per 1 \\ 3\per 2+1\per 0+2\per 3 & 3\per 1+1\per 4+2\per 1 \end{pmatrix}
=\begin{pmatrix} 2 & 9 \\ 12 & 9 \end{pmatrix}.
\]
\end{exemple}

\subsection{Exercicis resolts}

\begin{exresolt}[producte $2\times 2$ amb tots els passos]
Calcula $\begin{pmatrix} 4 & 2 \\ 1 & 3 \end{pmatrix}\cdot\begin{pmatrix} 1 & 5 \\ 2 & 0 \end{pmatrix}$.
\solucio
Escrivim cada casella com a suma de productes:
\begin{itemize}[leftmargin=1.4em, itemsep=2pt]
  \item Fila 1 $\times$ columna 1: $4\per 1+2\per 2=4+4=8$.
  \item Fila 1 $\times$ columna 2: $4\per 5+2\per 0=20+0=20$.
  \item Fila 2 $\times$ columna 1: $1\per 1+3\per 2=1+6=7$.
  \item Fila 2 $\times$ columna 2: $1\per 5+3\per 0=5+0=5$.
\end{itemize}
\[
\begin{pmatrix} 4 & 2 \\ 1 & 3 \end{pmatrix}\cdot\begin{pmatrix} 1 & 5 \\ 2 & 0 \end{pmatrix}=\begin{pmatrix} 8 & 20 \\ 7 & 5 \end{pmatrix}.
\]
\resposta{$\begin{pmatrix} 8 & 20 \\ 7 & 5 \end{pmatrix}$.}
\end{exresolt}

\begin{exresolt}[un producte rectangular $2\times 3$ per $3\times 1$]
Calcula $\begin{pmatrix} 2 & 1 & 3 \\ 0 & 4 & 1 \end{pmatrix}\cdot\begin{pmatrix} 1 \\ 2 \\ 3 \end{pmatrix}$.
\solucio
El resultat serà $2\times 1$ (files de la primera, columnes de la segona):
\begin{itemize}[leftmargin=1.4em, itemsep=2pt]
  \item Fila 1 $\times$ columna: $2\per 1+1\per 2+3\per 3=2+2+9=13$.
  \item Fila 2 $\times$ columna: $0\per 1+4\per 2+1\per 3=0+8+3=11$.
\end{itemize}
\[
\begin{pmatrix} 2 & 1 & 3 \\ 0 & 4 & 1 \end{pmatrix}\cdot\begin{pmatrix} 1 \\ 2 \\ 3 \end{pmatrix}=\begin{pmatrix} 13 \\ 11 \end{pmatrix}.
\]
\resposta{$\begin{pmatrix} 13 \\ 11 \end{pmatrix}$.}
\end{exresolt}

\subsection{Exercicis per fer}

\textit{Bateria graduada: de més senzill (1) a més complet (7). Escriu tots els productes abans de sumar.}

\begin{exercicis}
\begin{exlist}
  \item $\begin{pmatrix} 1 & 0 \\ 0 & 1 \end{pmatrix}\cdot\begin{pmatrix} 7 & 3 \\ 2 & 5 \end{pmatrix}$
  \item $\begin{pmatrix} 2 & 1 \\ 3 & 0 \end{pmatrix}\cdot\begin{pmatrix} 4 & 2 \\ 1 & 5 \end{pmatrix}$
  \item $\begin{pmatrix} 3 & 2 \\ 1 & 4 \end{pmatrix}\cdot\begin{pmatrix} 2 & 0 \\ 3 & 1 \end{pmatrix}$
  \item $\begin{pmatrix} 5 & 1 \\ 2 & 3 \end{pmatrix}\cdot\begin{pmatrix} 1 & 2 \\ 4 & 0 \end{pmatrix}$
  \item $\begin{pmatrix} 1 & 2 & 1 \\ 0 & 3 & 2 \end{pmatrix}\cdot\begin{pmatrix} 2 \\ 1 \\ 3 \end{pmatrix}$
  \item $\begin{pmatrix} 2 & 0 & 1 \\ 1 & 3 & 2 \end{pmatrix}\cdot\begin{pmatrix} 1 & 2 \\ 0 & 1 \\ 4 & 1 \end{pmatrix}$
  \item \textbf{(Raonament)} Calcula $A\cdot B$ i $B\cdot A$ amb
        $A=\begin{pmatrix} 1 & 2 \\ 0 & 1 \end{pmatrix}$ i
        $B=\begin{pmatrix} 3 & 0 \\ 1 & 2 \end{pmatrix}$. Surten iguals? Què en
        conclous sobre l'ordre dels factors?
\end{exlist}
\end{exercicis}

% ---------------------------------------------------------------------
\fullsolucions{Sessió 9}
\begin{solucions}
\begin{sollist}
  \item $\begin{pmatrix} 7 & 3 \\ 2 & 5 \end{pmatrix}$ (multiplicar per la matriu
        identitat deixa la matriu igual).
  \item $\begin{pmatrix} 9 & 9 \\ 12 & 6 \end{pmatrix}$.
  \item $\begin{pmatrix} 12 & 2 \\ 14 & 4 \end{pmatrix}$.
  \item $\begin{pmatrix} 9 & 10 \\ 14 & 4 \end{pmatrix}$.
  \item $\begin{pmatrix} 7 \\ 9 \end{pmatrix}$.
  \item $\begin{pmatrix} 6 & 5 \\ 9 & 7 \end{pmatrix}$.
  \item $A\cdot B=\begin{pmatrix} 5 & 4 \\ 1 & 2 \end{pmatrix}$,
        \;$B\cdot A=\begin{pmatrix} 3 & 6 \\ 1 & 4 \end{pmatrix}$. No surten iguals:
        en el producte de matrices l'ordre dels factors importa (en general
        $A\cdot B\neq B\cdot A$).
\end{sollist}
\end{solucions}
